% Abstract (English)
\newpage
\phantomsection
\addcontentsline{toc}{matter}{ABSTRACT}
\begin{spacing}{1.15}
\begin{center}
    \textbf{ABSTRACT}
\end{center}

Artificial intelligence (AI) models such as DeepBach and Coconet can harmonize a soprano melody in four parts, and their output has reached non-specialist users. In the Bach Doodle analysis, Huang et al.\ (2019) counted parallel fifths and octaves in Coconet harmonizations automatically. That finding came from uncontrolled usage data, one model, and two rules. This quantitative study continues the analysis under controlled conditions. The same melodies, Bach chorale sopranos, are given to DeepBach and Coconet. The context is limited to chorale harmonization, which in Gustav Strube's harmony textbook uses Bach's melodies and treats the fermata as a cadence. Rule adherence is measured as violations per measure for five voice-leading rules that belong to the Yan et al.\ (2018) rubric, are formulated after Strube, and can be decided without chord interpretation: parallel fifths, parallel octaves, upper-voice spacing, crossing above the soprano, and voice overlap. Bach's own harmonizations of the same melodies serve as a reference. All measurements are automatic, using the music21 library, without human raters. The instrument is checked against Strube's printed examples, music21's voice-leading functions, and the per-measure figures for Bach chorales reported by Huang et al. Differences between models, and between violations near fermatas and inside phrases for each model, are tested with the Wilcoxon signed-rank test. The study is expected to show which rules the models break most often and where in the phrase the violations occur, compared with Bach's harmonizations. Findings apply to the tested checkpoints and are not intended to establish the superiority of any model architecture.\par

\noindent\textbf{Keywords:} Automatic Harmonization, Voice-Leading Rules, Gustav Strube, Chorale, DeepBach, Coconet
\end{spacing}
